\section{Impatto sulle performance}

Vantaggi di Code Splitting e Lazy Loading:
\begin{itemize}
    \item Riduzione dei tempi di caricamento iniziale: Code splitting divide il codice JavaScript in moduli che vengono caricati solo quando necessari, permettendo alle pagine di caricarsi più velocemente. Lazy loading, invece, posticipa il caricamento di risorse (come immagini o video) fino a quando non sono effettivamente visualizzati, riducendo la quantità di dati inizialmente scaricati
    \item Migliore esperienza utente: I tempi di caricamento più rapidi e la maggiore reattività dell'applicazione consentono agli utenti di iniziare a interagire con il sito prima che tutte le risorse siano caricate. Questo è particolarmente utile per siti con molti contenuti multimediali o applicazioni JavaScript complesse
    \item Riduzione dell'uso di banda: Lazy loading carica solo le risorse necessarie, ottimizzando il consumo di banda, specialmente su dispositivi mobili e connessioni lente
    \item Migliore gestione delle risorse: Queste tecniche riducono il carico sui server e migliorano la gestione delle risorse di memoria, poiché vengono gestite in modo più efficiente le richieste di risorse non immediatamente necessarie
\end{itemize}

Sfide e svantaggi:
\begin{itemize}
    \item LCP peggiorato: Lazy loading può peggiorare il Largest Contentful Paint (LCP) su alcune pagine, poiché ritarda il caricamento delle immagini necessarie nella parte visibile (above-the-fold). Studi su WordPress hanno mostrato che disabilitare il lazy loading per le immagini sopra la piega può migliorare significativamente il LCP
    \item Gestione complessa: Implementare correttamente code splitting e lazy loading richiede una comprensione approfondita di quando e come dividere il codice o posticipare il caricamento delle risorse. Configurazioni errate possono portare a problemi di performance o alla mancata visualizzazione di contenuti essenziali
    \item Possibili conflitti con altre ottimizzazioni: Code splitting e lazy loading devono essere ben bilanciati con altre tecniche come il prefetching e il caching, per evitare conflitti di banda che potrebbero rallentare l'esperienza utente
\end{itemize}

\paragraph{Studi sull'impatto delle performance}
Alcuni studi dimostrano l'importanza di implementare correttamente queste tecniche. Ad esempio, un test su WordPress ha rilevato che disabilitare il lazy loading migliorava il LCP su pagine di archivio del 13-15\%, mentre l'abilitazione del lazy loading riduceva la quantità di dati immagine caricati, migliorando la banda utilizzata, ma a scapito di un ritardo nel caricamento delle immagini principali

Nel complesso, quando usate correttamente, queste tecniche forniscono un ottimo equilibrio tra prestazioni e esperienza utente, ma richiedono un'implementazione oculata per evitare effetti indesiderati.
